\documentclass[11pt,a4paper]{article}

\usepackage[utf8]{inputenc}
\usepackage[T1]{fontenc}
\usepackage[french]{babel}
\usepackage[a4paper,margin=2.2cm]{geometry}
\usepackage{amsmath,amssymb,amsthm,mathtools}
\usepackage{lmodern}
\usepackage{enumitem}
\usepackage{xcolor}
\usepackage{hyperref}
\usepackage{fancyhdr}
\usepackage{stmaryrd}
\usepackage{tikz}
\usepackage{mathrsfs}

\hypersetup{
    colorlinks=true,
    linkcolor=black,
    urlcolor=blue
}

\setlength{\parindent}{0pt}
\setlength{\parskip}{0.5em}

\pagestyle{fancy}
\fancyhf{}
\fancyhead[L]{Nombres complexes --- Exercices de colles}
\fancyhead[R]{Terminale $\rightarrow$ Prépa}
\fancyfoot[C]{\thepage}

\newtheorem*{definition}{Définition}
\newtheorem*{remarque}{Remarque}
\newtheorem*{methode}{Méthode}
\newtheorem*{idee}{Idée}
\newtheorem*{corrige}{Corrigé}
\newtheorem*{bilan}{Bilan}

\newcommand{\C}{\mathbb{C}}
\newcommand{\R}{\mathbb{R}}
\newcommand{\ii}{\mathrm{i}}
\newcommand{\e}{\mathrm{e}}
\newcommand{\Arg}{\mathrm{Arg}}
\newcommand{\Ree}{\mathrm{Re}}
\newcommand{\Imm}{\mathrm{Im}}

% Étoiles compatibles LaTeX classique
\newcommand{\fullstar}{\(\star\)}
\newcommand{\emptystar}{\(\circ\)}

\newcommand{\oneStar}{\fullstar\emptystar\emptystar\emptystar\emptystar}
\newcommand{\twoStars}{\fullstar\fullstar\emptystar\emptystar\emptystar}
\newcommand{\threeStars}{\fullstar\fullstar\fullstar\emptystar\emptystar}
\newcommand{\fourStars}{\fullstar\fullstar\fullstar\fullstar\emptystar}
\newcommand{\fiveStars}{\fullstar\fullstar\fullstar\fullstar\fullstar}

\newcommand{\difficulty}[1]{\textbf{Difficulté : #1}}

\title{\Huge Exercices de colles sur les nombres complexes\\[0.5em]
\Large Niveau Terminale vers Prépa\\[0.3em]
\large Avec corrigés détaillés}
\author{}
\date{}

\begin{document}

\maketitle

\tableofcontents
\newpage

\section*{Introduction}
Ce document propose une série d'exercices de colles sur les nombres complexes, classés implicitement du plus accessible au plus exigeant, pour un élève de Terminale souhaitant préparer la transition vers la prépa.

Les objectifs sont les suivants :
\begin{itemize}[label=--]
    \item maîtriser les calculs algébriques de base sur les complexes ;
    \item comprendre le rôle du module, du conjugué et de l'argument ;
    \item savoir passer de la forme algébrique à la forme trigonométrique ;
    \item commencer à faire de la géométrie avec les complexes ;
    \item développer des raisonnements plus structurés, de type prépa.
\end{itemize}

\medskip

Échelle proposée :
\begin{itemize}[label=--]
    \item \oneStar\ : bases Terminale ;
    \item \twoStars\ : Terminale solide ;
    \item \threeStars\ : transition vers la prépa ;
    \item \fourStars\ : bon niveau de début de prépa ;
    \item \fiveStars\ : exercice déjà assez fin, excellent pour une très bonne transition.
\end{itemize}

\newpage

% ============================================================
\section{Exercice 1 \oneStar\ --- Calculs de base}
\difficulty{\oneStar}

On considère
\[
z_1=3-2\ii, \qquad z_2=1+4\ii.
\]

\begin{enumerate}
    \item Calculer \(z_1+z_2\), \(z_1-z_2\), \(z_1z_2\).
    \item Calculer \(\dfrac{z_1}{z_2}\) sous forme algébrique.
    \item Déterminer le module et un argument de \(z_1\).
    \item Déterminer le conjugué de \(\dfrac{z_1}{z_2}\).
\end{enumerate}

\begin{corrige}
\textbf{1) Somme, différence, produit}

On calcule directement :
\[
z_1+z_2=(3-2\ii)+(1+4\ii)=4+2\ii.
\]

\[
z_1-z_2=(3-2\ii)-(1+4\ii)=2-6\ii.
\]

Pour le produit :
\[
z_1z_2=(3-2\ii)(1+4\ii)=3+12\ii-2\ii-8\ii^2.
\]
Or \(\ii^2=-1\), donc
\[
-8\ii^2=8.
\]
Ainsi
\[
z_1z_2=3+10\ii+8=11+10\ii.
\]

\textbf{2) Quotient}

On écrit
\[
\frac{z_1}{z_2}=\frac{3-2\ii}{1+4\ii}.
\]
Pour mettre sous forme algébrique, on multiplie numérateur et dénominateur par le conjugué du dénominateur, soit \(1-4\ii\) :
\[
\frac{3-2\ii}{1+4\ii}\times \frac{1-4\ii}{1-4\ii}
= \frac{(3-2\ii)(1-4\ii)}{(1+4\ii)(1-4\ii)}.
\]

Calculons le numérateur :
\[
(3-2\ii)(1-4\ii)=3-12\ii-2\ii+8\ii^2=3-14\ii-8=-5-14\ii.
\]

Calculons le dénominateur :
\[
(1+4\ii)(1-4\ii)=1-(4\ii)^2=1-16\ii^2=1+16=17.
\]

Donc
\[
\frac{z_1}{z_2}=\frac{-5-14\ii}{17}=-\frac{5}{17}-\frac{14}{17}\ii.
\]

\textbf{3) Module et argument de \(z_1\)}

On a
\[
z_1=3-2\ii.
\]
Son module vaut
\[
|z_1|=\sqrt{3^2+(-2)^2}=\sqrt{9+4}=\sqrt{13}.
\]

Pour l'argument, on regarde la position du point \((3,-2)\) dans le plan : il est dans le quatrième quadrant. On peut donc prendre comme argument
\[
\theta=-\arctan\left(\frac{2}{3}\right).
\]

Ainsi,
\[
|z_1|=\sqrt{13},
\qquad
\arg(z_1)=-\arctan\left(\frac{2}{3}\right)\ \text{mod }2\pi.
\]

\textbf{4) Conjugué du quotient}

On a trouvé
\[
\frac{z_1}{z_2}=-\frac{5}{17}-\frac{14}{17}\ii.
\]
Son conjugué est donc
\[
\overline{\left(\frac{z_1}{z_2}\right)}=-\frac{5}{17}+\frac{14}{17}\ii.
\]

\textbf{Conclusion :}
\[
z_1+z_2=4+2\ii,\qquad
z_1-z_2=2-6\ii,\qquad
z_1z_2=11+10\ii,
\]
\[
\frac{z_1}{z_2}=-\frac{5}{17}-\frac{14}{17}\ii,
\qquad
|z_1|=\sqrt{13},
\qquad
\arg(z_1)=-\arctan\left(\frac{2}{3}\right),
\]
\[
\overline{\left(\frac{z_1}{z_2}\right)}=-\frac{5}{17}+\frac{14}{17}\ii.
\]
\end{corrige}

% ============================================================
\section{Exercice 2 \oneStar\ --- Équation du second degré}
\difficulty{\oneStar}

Résoudre dans \(\C\) :
\[
z^2-6z+13=0.
\]

Puis écrire les solutions sous forme algébrique et donner leur module.

\begin{corrige}
On résout l'équation comme une équation du second degré classique.

\[
z^2-6z+13=0.
\]

Le discriminant vaut
\[
\Delta=(-6)^2-4\times 1\times 13=36-52=-16.
\]

Dans \(\C\), on peut écrire
\[
\sqrt{\Delta}=\sqrt{-16}=4\ii.
\]

Ainsi les solutions sont
\[
z=\frac{6\pm 4\ii}{2}=3\pm 2\ii.
\]

Les deux solutions sont donc :
\[
z_1=3+2\ii,\qquad z_2=3-2\ii.
\]

\textbf{Module des solutions}

Pour \(z_1=3+2\ii\),
\[
|z_1|=\sqrt{3^2+2^2}=\sqrt{13}.
\]

Pour \(z_2=3-2\ii\),
\[
|z_2|=\sqrt{3^2+(-2)^2}=\sqrt{13}.
\]

\textbf{Conclusion :}
\[
\boxed{z=3+2\ii \quad \text{ou} \quad z=3-2\ii}
\]
et dans les deux cas
\[
|z|=\sqrt{13}.
\]
\end{corrige}

% ============================================================
\section{Exercice 3 \oneStar\ --- Lieux géométriques simples}
\difficulty{\oneStar}

\begin{enumerate}
    \item Déterminer l'ensemble des points \(M\) d'affixe \(z\) tels que
    \[
    |z-2|=3.
    \]
    \item Interpréter géométriquement l'ensemble des points vérifiant
    \[
    |z-1|=|z+3|.
    \]
\end{enumerate}

\begin{corrige}
\textbf{1) Étude de \(|z-2|=3\)}

Soit \(M\) le point d'affixe \(z\), et \(A\) le point d'affixe \(2\). Alors
\[
|z-2|=MA.
\]
La condition
\[
|z-2|=3
\]
signifie donc que la distance de \(M\) à \(A\) vaut \(3\).

L'ensemble cherché est donc le cercle de centre \(A\) d'affixe \(2\) et de rayon \(3\).

\[
\boxed{\text{Cercle de centre }2\text{ et de rayon }3}
\]

\textbf{2) Étude de \(|z-1|=|z+3|\)}

Soient \(A\) d'affixe \(1\) et \(B\) d'affixe \(-3\). Alors
\[
|z-1|=MA,\qquad |z+3|=MB.
\]
La condition
\[
|z-1|=|z+3|
\]
signifie que \(M\) est à égale distance de \(A\) et de \(B\).

L'ensemble des points équidistants de deux points fixes est la médiatrice du segment \([AB]\).

Comme \(A\) et \(B\) sont sur l'axe réel, la médiatrice est la droite verticale passant par le milieu de \([AB]\). Le milieu a pour abscisse
\[
\frac{1+(-3)}{2}=-1.
\]
Donc l'ensemble est la droite d'équation
\[
x=-1.
\]

\[
\boxed{\text{Médiatrice de }[AB],\ \text{soit la droite }x=-1}
\]
\end{corrige}

% ============================================================
\section{Exercice 4 \oneStar\ --- Argument d'un quotient}
\difficulty{\oneStar}

On considère
\[
z=\frac{1+\ii}{1-\ii}.
\]

\begin{enumerate}
    \item Mettre \(z\) sous forme algébrique.
    \item Déterminer son module et un argument.
    \item En déduire une valeur de \(\arg\left(\dfrac{1+\ii}{1-\ii}\right)\).
\end{enumerate}

\begin{corrige}
\textbf{1) Forme algébrique}

On rationalise :
\[
z=\frac{1+\ii}{1-\ii}\times \frac{1+\ii}{1+\ii}
=\frac{(1+\ii)^2}{1-\ii^2}.
\]

Calcul du numérateur :
\[
(1+\ii)^2=1+2\ii+\ii^2=1+2\ii-1=2\ii.
\]

Calcul du dénominateur :
\[
1-\ii^2=1-(-1)=2.
\]

Donc
\[
z=\frac{2\ii}{2}=\ii.
\]

\textbf{2) Module et argument}

Comme \(z=\ii\), on a
\[
|z|=1.
\]
Un argument de \(\ii\) est
\[
\frac{\pi}{2}.
\]

\textbf{3) Conclusion}

Ainsi
\[
\arg\left(\frac{1+\ii}{1-\ii}\right)=\frac{\pi}{2}\quad \text{mod }2\pi.
\]

\[
\boxed{\frac{1+\ii}{1-\ii}=\ii,\qquad |z|=1,\qquad \arg(z)=\frac{\pi}{2}}
\]
\end{corrige}

% ============================================================
\section{Exercice 5 \twoStars\ --- Forme trigonométrique et puissances}
\difficulty{\twoStars}

On pose
\[
z=\frac{\sqrt3}{2}+\frac12 \ii.
\]

\begin{enumerate}
    \item Donner le module et un argument de \(z\).
    \item Écrire \(z\) sous forme trigonométrique.
    \item Calculer \(z^n\) pour tout entier naturel \(n\).
    \item Déterminer les entiers \(n\) tels que \(z^n\in\R\).
\end{enumerate}

\begin{corrige}
\textbf{1) Module et argument}

On reconnaît
\[
\cos\left(\frac{\pi}{6}\right)=\frac{\sqrt3}{2},
\qquad
\sin\left(\frac{\pi}{6}\right)=\frac12.
\]
Ainsi
\[
z=\cos\left(\frac{\pi}{6}\right)+\ii\sin\left(\frac{\pi}{6}\right).
\]

Son module vaut
\[
|z|=\sqrt{\left(\frac{\sqrt3}{2}\right)^2+\left(\frac12\right)^2}
=\sqrt{\frac34+\frac14}=\sqrt1=1.
\]

Un argument de \(z\) est donc
\[
\frac{\pi}{6}.
\]

\textbf{2) Forme trigonométrique}

On peut écrire
\[
z=\cos\left(\frac{\pi}{6}\right)+\ii\sin\left(\frac{\pi}{6}\right).
\]
Ou encore
\[
z=\e^{\ii\pi/6}.
\]

\textbf{3) Calcul de \(z^n\)}

Par la formule de Moivre :
\[
z^n=\left(\cos\left(\frac{\pi}{6}\right)+\ii\sin\left(\frac{\pi}{6}\right)\right)^n
=\cos\left(\frac{n\pi}{6}\right)+\ii\sin\left(\frac{n\pi}{6}\right).
\]

Donc
\[
\boxed{z^n=\cos\left(\frac{n\pi}{6}\right)+\ii\sin\left(\frac{n\pi}{6}\right)}.
\]

\textbf{4) Condition pour que \(z^n\) soit réel}

Le nombre \(z^n\) est réel si et seulement si sa partie imaginaire est nulle, donc si
\[
\sin\left(\frac{n\pi}{6}\right)=0.
\]
Cela équivaut à
\[
\frac{n\pi}{6}=k\pi
\]
avec \(k\in\mathbb{Z}\), soit
\[
n=6k.
\]

Donc
\[
\boxed{z^n\in\R \iff 6\mid n}.
\]
\end{corrige}

% ============================================================
\section{Exercice 6 \twoStars\ --- Racines cubiques}
\difficulty{\twoStars}

Résoudre dans \(\C\) :
\[
z^3=8\ii.
\]

Donner les solutions sous forme trigonométrique puis algébrique.

\begin{corrige}
On écrit d'abord \(8\ii\) sous forme trigonométrique.

On a
\[
|8\ii|=8
\]
et un argument de \(8\ii\) est
\[
\frac{\pi}{2}.
\]
Donc
\[
8\ii=8\e^{\ii\pi/2}.
\]

On cherche les racines cubiques. Si
\[
z^3=8\e^{\ii\pi/2},
\]
alors les racines sont données par
\[
z_k=8^{1/3}\e^{\ii\left(\frac{\pi/2+2k\pi}{3}\right)}
=2\e^{\ii\left(\frac{\pi}{6}+\frac{2k\pi}{3}\right)},
\qquad k\in\{0,1,2\}.
\]

Donc les trois solutions sont :
\[
z_0=2\e^{\ii\pi/6},\qquad
z_1=2\e^{\ii5\pi/6},\qquad
z_2=2\e^{\ii3\pi/2}.
\]

\textbf{Passage à la forme algébrique}

Pour \(z_0\) :
\[
z_0=2\left(\cos\frac{\pi}{6}+\ii\sin\frac{\pi}{6}\right)
=2\left(\frac{\sqrt3}{2}+\frac12\ii\right)=\sqrt3+\ii.
\]

Pour \(z_1\) :
\[
z_1=2\left(\cos\frac{5\pi}{6}+\ii\sin\frac{5\pi}{6}\right)
=2\left(-\frac{\sqrt3}{2}+\frac12\ii\right)=-\sqrt3+\ii.
\]

Pour \(z_2\) :
\[
z_2=2\left(\cos\frac{3\pi}{2}+\ii\sin\frac{3\pi}{2}\right)=2(0-\ii)=-2\ii.
\]

\textbf{Conclusion :}
\[
\boxed{z\in\left\{\sqrt3+\ii,\ -\sqrt3+\ii,\ -2\ii\right\}}
\]
ou sous forme trigonométrique
\[
\boxed{z_k=2\e^{\ii\left(\frac{\pi}{6}+\frac{2k\pi}{3}\right)},\quad k=0,1,2.}
\]
\end{corrige}

% ============================================================
\section{Exercice 7 \twoStars\ --- Quotient de complexes et puissances}
\difficulty{\twoStars}

Soit
\[
z=\frac{2+\ii\sqrt3}{2-\ii\sqrt3}.
\]

\begin{enumerate}
    \item Montrer que \(|z|=1\).
    \item Déterminer un argument de \(z\).
    \item En déduire la forme trigonométrique de \(z\).
    \item Calculer \(z^{2026}\).
\end{enumerate}

\begin{corrige}
\textbf{1) Calcul du module}

On utilise la propriété
\[
\left|\frac{a}{b}\right|=\frac{|a|}{|b|}.
\]
Donc
\[
|z|=\frac{|2+\ii\sqrt3|}{|2-\ii\sqrt3|}.
\]
Or
\[
|2+\ii\sqrt3|=\sqrt{2^2+(\sqrt3)^2}=\sqrt{7},
\]
et
\[
|2-\ii\sqrt3|=\sqrt{2^2+(-\sqrt3)^2}=\sqrt{7}.
\]
Donc
\[
|z|=1.
\]

\textbf{2) Argument}

On sait que l'argument d'un quotient est la différence des arguments :
\[
\arg(z)=\arg(2+\ii\sqrt3)-\arg(2-\ii\sqrt3).
\]

Posons
\[
\theta=\arctan\left(\frac{\sqrt3}{2}\right).
\]
Alors
\[
\arg(2+\ii\sqrt3)=\theta,\qquad \arg(2-\ii\sqrt3)=-\theta.
\]
Donc
\[
\arg(z)=\theta-(-\theta)=2\theta.
\]

Ainsi
\[
\boxed{\arg(z)=2\arctan\left(\frac{\sqrt3}{2}\right)}.
\]

\textbf{3) Forme trigonométrique}

Comme \(|z|=1\), on peut écrire
\[
z=\e^{2\ii\theta}
\qquad \text{avec } \theta=\arctan\left(\frac{\sqrt3}{2}\right).
\]
Donc
\[
\boxed{z=\cos(2\theta)+\ii\sin(2\theta)}.
\]

\textbf{4) Calcul de \(z^{2026}\)}

On a
\[
z^{2026}=\e^{4052\ii\theta}.
\]
Donc
\[
\boxed{z^{2026}=\cos(4052\theta)+\ii\sin(4052\theta)}
\qquad \text{avec } \theta=\arctan\left(\frac{\sqrt3}{2}\right).
\]

\textbf{Remarque :} ici l'angle n'est pas un angle remarquable usuel, donc l'expression trigonométrique est déjà une forme correcte.
\end{corrige}

% ============================================================
\section{Exercice 8 \twoStars\ --- Partie réelle et module}
\difficulty{\twoStars}

Déterminer les nombres complexes \(z\) tels que
\[
z+\overline z=6
\qquad \text{et} \qquad
z\overline z=13.
\]

\begin{corrige}
Posons
\[
z=x+\ii y
\qquad \text{avec } x,y\in\R.
\]
Alors
\[
\overline z=x-\ii y.
\]

\textbf{Première condition}

\[
z+\overline z=(x+\ii y)+(x-\ii y)=2x.
\]
La condition
\[
z+\overline z=6
\]
donne donc
\[
2x=6 \quad \Longrightarrow \quad x=3.
\]

\textbf{Deuxième condition}

On sait que
\[
z\overline z=|z|^2=x^2+y^2.
\]
La condition
\[
z\overline z=13
\]
donne
\[
x^2+y^2=13.
\]
Or \(x=3\), donc
\[
9+y^2=13 \quad \Longrightarrow \quad y^2=4 \quad \Longrightarrow \quad y=\pm 2.
\]

Donc
\[
z=3+2\ii \quad \text{ou} \quad z=3-2\ii.
\]

\textbf{Conclusion :}
\[
\boxed{z=3+2\ii \quad \text{ou} \quad z=3-2\ii}
\]
\end{corrige}

% ============================================================
\section{Exercice 9 \threeStars\ --- Lieu géométrique}
\difficulty{\threeStars}

\begin{enumerate}
    \item Résoudre dans \(\C\) :
    \[
    |z-1|=|z-\ii|.
    \]
    \item Puis résoudre :
    \[
    |z-2|=2|z+1|.
    \]
\end{enumerate}

\begin{corrige}
\textbf{1) Étude de \(|z-1|=|z-\ii|\)}

Soit \(M\) d'affixe \(z\), \(A\) d'affixe \(1\), \(B\) d'affixe \(\ii\). Alors
\[
|z-1|=MA,\qquad |z-\ii|=MB.
\]
La condition signifie que
\[
MA=MB.
\]
L'ensemble des points équidistants de \(A\) et de \(B\) est la médiatrice du segment \([AB]\).

\[
\boxed{\text{Ensemble cherché : la médiatrice de }[AB]}
\]

\textbf{2) Étude de \(|z-2|=2|z+1|\)}

Posons \(z=x+\ii y\). Alors
\[
|z-2|=\sqrt{(x-2)^2+y^2},
\qquad
|z+1|=\sqrt{(x+1)^2+y^2}.
\]
L'équation devient
\[
\sqrt{(x-2)^2+y^2}=2\sqrt{(x+1)^2+y^2}.
\]
On élève au carré :
\[
(x-2)^2+y^2=4\big((x+1)^2+y^2\big).
\]
Développons :
\[
x^2-4x+4+y^2=4x^2+8x+4+4y^2.
\]
En ramenant tout à gauche :
\[
x^2-4x+4+y^2-4x^2-8x-4-4y^2=0,
\]
soit
\[
-3x^2-12x-3y^2=0.
\]
En divisant par \(-3\) :
\[
x^2+4x+y^2=0.
\]
On complète le carré :
\[
x^2+4x+4+y^2=4,
\]
donc
\[
(x+2)^2+y^2=4.
\]

C'est le cercle de centre \((-2,0)\) et de rayon \(2\).

\[
\boxed{(x+2)^2+y^2=4}
\]
\[
\boxed{\text{Cercle de centre }-2\text{ et de rayon }2}
\]
\end{corrige}

% ============================================================
\section{Exercice 10 \threeStars\ --- Triangle équilatéral et quotient}
\difficulty{\threeStars}

On considère trois points \(A,B,C\) d'affixes \(a,b,c\). Montrer que le triangle \(ABC\) est équilatéral direct si et seulement si
\[
\frac{c-a}{b-a}=e^{\ii\pi/3}.
\]

Puis écrire une condition analogue pour un triangle équilatéral indirect.

\begin{corrige}
On suppose d'abord \(a\neq b\), sinon le triangle serait dégénéré.

\textbf{Interprétation du quotient}

Le quotient
\[
\frac{c-a}{b-a}
\]
compare le vecteur \(\overrightarrow{AC}\) au vecteur \(\overrightarrow{AB}\).

En effet :
\begin{itemize}
    \item son module vaut
    \[
    \left|\frac{c-a}{b-a}\right|=\frac{|c-a|}{|b-a|}=\frac{AC}{AB},
    \]
    \item son argument vaut
    \[
    \arg\left(\frac{c-a}{b-a}\right)=\arg(c-a)-\arg(b-a),
    \]
    ce qui est précisément l'angle orienté entre \(\overrightarrow{AB}\) et \(\overrightarrow{AC}\).
\end{itemize}

\textbf{Sens direct}

Si le triangle \(ABC\) est équilatéral direct, alors :
\begin{itemize}
    \item \(AB=AC\), donc
    \[
    \left|\frac{c-a}{b-a}\right|=1,
    \]
    \item l'angle orienté de \(\overrightarrow{AB}\) vers \(\overrightarrow{AC}\) vaut \(\dfrac{\pi}{3}\), donc
    \[
    \arg\left(\frac{c-a}{b-a}\right)=\frac{\pi}{3}\ \text{mod }2\pi.
    \]
\end{itemize}
Un complexe de module \(1\) et d'argument \(\dfrac{\pi}{3}\) est précisément
\[
e^{\ii\pi/3}.
\]
Donc
\[
\frac{c-a}{b-a}=e^{\ii\pi/3}.
\]

\textbf{Réciproque}

Supposons
\[
\frac{c-a}{b-a}=e^{\ii\pi/3}.
\]
Alors le module du quotient vaut \(1\), donc
\[
|c-a|=|b-a|,
\]
c'est-à-dire
\[
AC=AB.
\]
De plus l'argument du quotient vaut \(\dfrac{\pi}{3}\), donc l'angle orienté entre \(\overrightarrow{AB}\) et \(\overrightarrow{AC}\) vaut \(\dfrac{\pi}{3}\). Le triangle est donc équilatéral direct.

\textbf{Cas indirect}

Pour un triangle équilatéral indirect, l'angle orienté vaut
\[
-\frac{\pi}{3}.
\]
La condition devient alors
\[
\boxed{\frac{c-a}{b-a}=e^{-\ii\pi/3}}.
\]

\textbf{Conclusion :}
\[
\boxed{ABC\ \text{équilatéral direct} \iff \frac{c-a}{b-a}=e^{\ii\pi/3}}
\]
et
\[
\boxed{ABC\ \text{équilatéral indirect} \iff \frac{c-a}{b-a}=e^{-\ii\pi/3}}.
\]
\end{corrige}

% ============================================================
\section{Exercice 11 \threeStars\ --- Somme géométrique complexe}
\difficulty{\threeStars}

Soit \(z\neq 1\). Calculer
\[
S_n=1+z+z^2+\cdots+z^n.
\]

Application : calculer
\[
1+e^{\ii\theta}+e^{2\ii\theta}+\cdots+e^{n\ii\theta}.
\]

Puis en déduire une expression de
\[
\sum_{k=0}^n \cos(k\theta)
\qquad \text{et} \qquad
\sum_{k=0}^n \sin(k\theta).
\]

\begin{corrige}
\textbf{1) Somme géométrique}

On pose
\[
S_n=1+z+z^2+\cdots+z^n.
\]
On multiplie par \(z\) :
\[
zS_n=z+z^2+\cdots+z^{n+1}.
\]
En soustrayant :
\[
S_n-zS_n=1-z^{n+1}.
\]
Donc
\[
(1-z)S_n=1-z^{n+1}.
\]
Comme \(z\neq 1\), on peut diviser :
\[
\boxed{S_n=\frac{1-z^{n+1}}{1-z}}.
\]

\textbf{2) Application à \(z=e^{\ii\theta}\)}

Si \(e^{\ii\theta}\neq 1\), alors
\[
1+e^{\ii\theta}+e^{2\ii\theta}+\cdots+e^{n\ii\theta}
=\frac{1-e^{(n+1)\ii\theta}}{1-e^{\ii\theta}}.
\]

\textbf{3) Sommes des cosinus et des sinus}

Or
\[
e^{k\ii\theta}=\cos(k\theta)+\ii\sin(k\theta).
\]
Donc
\[
\sum_{k=0}^n e^{k\ii\theta}
=
\sum_{k=0}^n \cos(k\theta)
+\ii \sum_{k=0}^n \sin(k\theta).
\]
La partie réelle donne
\[
\sum_{k=0}^n \cos(k\theta)=\Ree\left(\frac{1-e^{(n+1)\ii\theta}}{1-e^{\ii\theta}}\right),
\]
et la partie imaginaire donne
\[
\sum_{k=0}^n \sin(k\theta)=\Imm\left(\frac{1-e^{(n+1)\ii\theta}}{1-e^{\ii\theta}}\right).
\]

\textbf{Forme plus classique}

On peut aussi obtenir
\[
\sum_{k=0}^n \cos(k\theta)
=
\frac{\sin\left(\frac{(n+1)\theta}{2}\right)\cos\left(\frac{n\theta}{2}\right)}
{\sin\left(\frac{\theta}{2}\right)},
\]
et
\[
\sum_{k=0}^n \sin(k\theta)
=
\frac{\sin\left(\frac{(n+1)\theta}{2}\right)\sin\left(\frac{n\theta}{2}\right)}
{\sin\left(\frac{\theta}{2}\right)},
\]
lorsque \(\theta\notin 2\pi\mathbb{Z}\).

\textbf{Conclusion principale :}
\[
\boxed{1+z+\cdots+z^n=\frac{1-z^{n+1}}{1-z}}
\]
pour \(z\neq 1\).
\end{corrige}

% ============================================================
\section{Exercice 12 \threeStars\ --- Racines de l'unité}
\difficulty{\threeStars}

Déterminer les solutions de
\[
z^5=1.
\]

\begin{enumerate}
    \item Les écrire explicitement.
    \item Les représenter dans le plan.
    \item Calculer leur somme.
    \item Calculer
    \[
    (1-z_1)(1-z_2)(1-z_3)(1-z_4)
    \]
    où \(z_1,z_2,z_3,z_4\) sont les racines cinquièmes de l'unité distinctes de \(1\).
\end{enumerate}

\begin{corrige}
\textbf{1) Solutions de \(z^5=1\)}

On écrit
\[
1=e^{2k\ii\pi}
\]
pour tout entier \(k\). Les racines cinquièmes sont donc
\[
z_k=e^{2k\ii\pi/5},\qquad k=0,1,2,3,4.
\]

Ainsi,
\[
\boxed{z_k=e^{2k\ii\pi/5}\quad (k=0,1,2,3,4)}.
\]

\textbf{2) Représentation géométrique}

Ces cinq racines sont réparties régulièrement sur le cercle unité. Elles forment les sommets d'un pentagone régulier inscrit dans le cercle trigonométrique.

\textbf{3) Somme des racines}

Les racines de \(z^5-1=0\) vérifient
\[
z^5-1=(z-1)(z^4+z^3+z^2+z+1).
\]
Donc si \(z\neq 1\) et \(z^5=1\), alors
\[
z^4+z^3+z^2+z+1=0.
\]
La somme des cinq racines vaut donc
\[
1+z_1+z_2+z_3+z_4=0.
\]

\[
\boxed{\sum_{k=0}^4 z_k=0}
\]

\textbf{4) Calcul du produit}

On a
\[
z^5-1=(z-1)(z-z_1)(z-z_2)(z-z_3)(z-z_4).
\]
On utilise
\[
z^4+z^3+z^2+z+1=(z-z_1)(z-z_2)(z-z_3)(z-z_4).
\]
En prenant \(z=1\), on obtient
\[
1+1+1+1+1=(1-z_1)(1-z_2)(1-z_3)(1-z_4).
\]
Donc
\[
\boxed{(1-z_1)(1-z_2)(1-z_3)(1-z_4)=5}.
\]
\end{corrige}

% ============================================================
\section{Exercice 13 \fourStars\ --- Quand \(z+\frac1z\) est réel}
\difficulty{\fourStars}

Déterminer tous les complexes \(z\) tels que
\[
z+\frac1z\in\R
\qquad \text{avec } z\neq 0.
\]

Montrer que cela équivaut à
\[
z\in\R \quad \text{ou} \quad |z|=1.
\]

\begin{corrige}
On veut que
\[
z+\frac1z
\]
soit réel.

Un complexe est réel si et seulement s'il est égal à son conjugué. Donc on impose
\[
z+\frac1z=\overline z+\frac1{\overline z}.
\]

On réécrit :
\[
z-\overline z+\frac1z-\frac1{\overline z}=0.
\]
Or
\[
\frac1z-\frac1{\overline z}=\frac{\overline z-z}{z\overline z}.
\]
Donc
\[
z-\overline z+\frac{\overline z-z}{z\overline z}=0.
\]
On factorise par \(z-\overline z\) :
\[
(z-\overline z)\left(1-\frac1{z\overline z}\right)=0.
\]

Ainsi :
\[
z-\overline z=0
\qquad \text{ou} \qquad
1-\frac1{|z|^2}=0.
\]

Le premier cas donne
\[
z=\overline z \iff z\in\R.
\]

Le second donne
\[
\frac1{|z|^2}=1 \iff |z|^2=1 \iff |z|=1.
\]

Donc
\[
\boxed{z+\frac1z\in\R \iff z\in\R \ \text{ou}\ |z|=1}
\]
pour \(z\neq 0\).

\textbf{Conclusion :}
les solutions sont exactement les complexes non nuls situés sur l'axe réel ou sur le cercle unité.
\end{corrige}

% ============================================================
\section{Exercice 14 \fourStars\ --- Étude d'une transformation}
\difficulty{\fourStars}

On considère l'application
\[
f(z)=\frac{z-1}{z+1},
\qquad z\neq -1.
\]

\begin{enumerate}
    \item Déterminer l'image des réels par \(f\).
    \item Déterminer les \(z\) tels que \(f(z)\) soit imaginaire pur.
    \item Déterminer les \(z\) tels que \(|f(z)|=1\).
\end{enumerate}

\begin{corrige}
\textbf{1) Image des réels}

Si \(z=x\in\R\) et \(x\neq -1\), alors
\[
f(x)=\frac{x-1}{x+1}.
\]
C'est un quotient de deux réels, donc c'est un réel.

Ainsi
\[
\boxed{f(\R\setminus\{-1\})\subset \R}.
\]

Réciproquement, si \(y\in\R\) et \(y\neq 1\), on peut résoudre
\[
y=\frac{x-1}{x+1}.
\]
Cela donne
\[
y(x+1)=x-1
\]
puis
\[
yx+y=x-1
\]
donc
\[
x(y-1)=-(1+y).
\]
Si \(y\neq 1\), alors
\[
x=\frac{-(1+y)}{y-1}=\frac{1+y}{1-y}\in\R.
\]
Donc les réels sont bien atteints sauf éventuellement \(1\), qu'on n'obtient pas.

Ainsi l'image des réels est
\[
\boxed{\R\setminus\{1\}}.
\]

\textbf{2) Quand \(f(z)\) est imaginaire pur}

Un complexe est imaginaire pur si sa partie réelle est nulle. Posons
\[
z=x+\ii y.
\]
Alors
\[
f(z)=\frac{x-1+\ii y}{x+1+\ii y}.
\]
On multiplie par le conjugué du dénominateur :
\[
f(z)=\frac{(x-1+\ii y)(x+1-\ii y)}{(x+1)^2+y^2}.
\]
Développons le numérateur :
\[
(x-1)(x+1)+y^2+\ii\big(y(x+1)-y(x-1)\big).
\]
La partie réelle vaut donc
\[
x^2-1+y^2.
\]
Pour que \(f(z)\) soit imaginaire pur, il faut et il suffit que
\[
x^2+y^2-1=0.
\]
Donc
\[
\boxed{|z|=1}
\]
avec \(z\neq -1\) puisque \(f\) n'est pas définie en \(-1\).

\textbf{3) Quand \(|f(z)|=1\)}

On calcule :
\[
\left|\frac{z-1}{z+1}\right|=1
\iff |z-1|=|z+1|.
\]
C'est l'ensemble des points équidistants de \(1\) et de \(-1\). Il s'agit donc de la médiatrice du segment joignant ces deux points, c'est-à-dire l'axe imaginaire.

Donc
\[
\boxed{|f(z)|=1 \iff \Ree(z)=0.}
\]
\end{corrige}

% ============================================================
\section{Exercice 15 \threeStars\ --- Racines carrées d'un complexe}
\difficulty{\threeStars}

Déterminer les solutions de
\[
z^2=5+12\ii.
\]

On cherchera les solutions sous la forme \(z=x+\ii y\), avec \(x,y\in\R\).

\begin{corrige}
Posons
\[
z=x+\ii y.
\]
Alors
\[
z^2=(x+\ii y)^2=x^2-y^2+2xy\ii.
\]
On veut
\[
x^2-y^2+2xy\ii=5+12\ii.
\]
On identifie partie réelle et partie imaginaire :
\[
\begin{cases}
x^2-y^2=5,\\
2xy=12.
\end{cases}
\]
La deuxième équation donne
\[
xy=6.
\]

On cherche deux réels dont le produit vaut \(6\) et dont la différence des carrés vaut \(5\).

Essayons les valeurs entières simples : \(x=3\), \(y=2\) donne
\[
xy=6,\qquad x^2-y^2=9-4=5.
\]
Donc
\[
3+2\ii
\]
est une solution.

L'autre racine carrée est l'opposée :
\[
-(3+2\ii)=-3-2\ii.
\]

Vérification :
\[
(-3-2\ii)^2=(3+2\ii)^2=5+12\ii.
\]

\textbf{Conclusion :}
\[
\boxed{z=3+2\ii \quad \text{ou} \quad z=-3-2\ii.}
\]
\end{corrige}

% ============================================================
\section{Exercice 16 \fourStars\ --- Équation bicarrée en apparence}
\difficulty{\fourStars}

Résoudre dans \(\C\) :
\[
z^4+(1-\ii)z^2-2\ii=0.
\]

\begin{corrige}
L'idée est de poser
\[
Z=z^2.
\]
L'équation devient
\[
Z^2+(1-\ii)Z-2\ii=0.
\]

Calculons le discriminant :
\[
\Delta=(1-\ii)^2-4(-2\ii).
\]
Or
\[
(1-\ii)^2=1-2\ii+\ii^2=1-2\ii-1=-2\ii.
\]
Donc
\[
\Delta=-2\ii+8\ii=6\ii.
\]

Il faut donc résoudre
\[
Z=\frac{-(1-\ii)\pm \sqrt{6\ii}}{2}.
\]

Cherchons \(\sqrt{6\ii}\). On écrit
\[
6\ii=6\e^{\ii\pi/2}.
\]
Ainsi
\[
\sqrt{6\ii}=\sqrt6\,\e^{\ii\pi/4}=\sqrt6\left(\frac{\sqrt2}{2}+\ii\frac{\sqrt2}{2}\right)=\sqrt3+\ii\sqrt3.
\]

Donc
\[
Z=\frac{-1+\ii\pm (\sqrt3+\ii\sqrt3)}{2}.
\]

On obtient deux valeurs :
\[
Z_1=\frac{-1+\sqrt3}{2}+\ii\frac{1+\sqrt3}{2},
\]
\[
Z_2=\frac{-1-\sqrt3}{2}+\ii\frac{1-\sqrt3}{2}.
\]

Il faut ensuite résoudre
\[
z^2=Z_1
\qquad \text{et} \qquad
z^2=Z_2.
\]

À ce stade, pour une colle de transition, on peut laisser les solutions sous la forme
\[
z=\pm \sqrt{Z_1},\qquad z=\pm \sqrt{Z_2}.
\]

\textbf{Conclusion :}
\[
\boxed{
z\in \left\{
\pm\sqrt{\frac{-1+\sqrt3}{2}+\ii\frac{1+\sqrt3}{2}},
\ 
\pm\sqrt{\frac{-1-\sqrt3}{2}+\ii\frac{1-\sqrt3}{2}}
\right\}.
}
\]

\textbf{Remarque :} le cœur de l'exercice est ici la substitution \(Z=z^2\), puis la résolution dans \(\C\).
\end{corrige}

% ============================================================
\section{Exercice 17 \fourStars\ --- Alignement et angle droit}
\difficulty{\fourStars}

Soient \(A,B,C\) d'affixes respectives \(a,b,c\), avec \(a\neq c\) et \(b\neq c\).

\begin{enumerate}
    \item Montrer que \(A,B,C\) sont alignés si et seulement si
    \[
    \frac{a-c}{b-c}\in\R.
    \]
    \item Montrer que \(A,B,C\) forment un angle droit en \(C\) si et seulement si
    \[
    \frac{a-c}{b-c}\in \ii\R.
    \]
\end{enumerate}

\begin{corrige}
\textbf{1) Alignement}

Les points \(A,B,C\) sont alignés si et seulement si les vecteurs \(\overrightarrow{CA}\) et \(\overrightarrow{CB}\) sont colinéaires.

Dans le langage des complexes, cela signifie que l'un est un multiple réel de l'autre :
\[
a-c=\lambda(b-c)
\qquad \text{avec } \lambda\in\R.
\]
Comme \(b\neq c\), cela équivaut à
\[
\frac{a-c}{b-c}=\lambda\in\R.
\]
Donc
\[
\boxed{A,B,C\ \text{sont alignés} \iff \frac{a-c}{b-c}\in\R.}
\]

\textbf{2) Angle droit en \(C\)}

Les vecteurs \(\overrightarrow{CA}\) et \(\overrightarrow{CB}\) sont perpendiculaires si et seulement si l'un s'obtient à partir de l'autre par multiplication par un réel puis par \(\ii\), c'est-à-dire par une rotation d'angle \(\pm \dfrac{\pi}{2}\).

Cela signifie qu'il existe \(\mu\in\R\) tel que
\[
a-c=\ii\mu (b-c).
\]
En divisant par \(b-c\), on obtient
\[
\frac{a-c}{b-c}=\ii\mu\in \ii\R.
\]

Réciproquement, si
\[
\frac{a-c}{b-c}\in \ii\R,
\]
alors le quotient a pour argument \(\pm \dfrac{\pi}{2}\) modulo \(\pi\), donc les vecteurs sont perpendiculaires.

Ainsi
\[
\boxed{A,B,C\ \text{forment un angle droit en }C \iff \frac{a-c}{b-c}\in \ii\R.}
\]
\end{corrige}

% ============================================================
\section{Exercice 18 \fiveStars\ --- Interprétation géométrique d'un argument}
\difficulty{\fiveStars}

Déterminer l'ensemble des points \(M\) d'affixe \(z\) tels que
\[
\arg\left(\frac{z-a}{z-b}\right)=\frac{\pi}{3},
\]
où \(a\) et \(b\) sont deux complexes fixés distincts.

\begin{corrige}
Soient \(A\) et \(B\) les points d'affixes respectives \(a\) et \(b\), et \(M\) le point d'affixe \(z\).

Le complexe
\[
\frac{z-a}{z-b}
\]
a pour argument
\[
\arg(z-a)-\arg(z-b).
\]
Or \(\arg(z-a)\) correspond à la direction du vecteur \(\overrightarrow{AM}\), et \(\arg(z-b)\) à la direction du vecteur \(\overrightarrow{BM}\).

Ainsi
\[
\arg\left(\frac{z-a}{z-b}\right)
\]
représente l'angle orienté entre les vecteurs \(\overrightarrow{BM}\) et \(\overrightarrow{AM}\), autrement dit l'angle \(\widehat{BMA}\) orienté.

La condition
\[
\arg\left(\frac{z-a}{z-b}\right)=\frac{\pi}{3}
\]
signifie donc que
\[
\widehat{BMA}=\frac{\pi}{3}.
\]

L'ensemble des points depuis lesquels le segment \([AB]\) est vu sous un angle constant \(\dfrac{\pi}{3}\) est un \emph{cercle capable}.

\textbf{Conclusion :}
\[
\boxed{\text{L'ensemble est l'un des deux arcs du cercle capable associé à }[AB]\text{ et à l'angle }\frac{\pi}{3}.}
\]

\textbf{Remarque :} le signe \(+\dfrac{\pi}{3}\) sélectionne un des deux arcs orientés possibles.
\end{corrige}

% ============================================================
\section{Exercice 19 \threeStars\ --- Polynôme et racines}
\difficulty{\threeStars}

On considère
\[
P(z)=z^3-3z+2.
\]

\begin{enumerate}
    \item Factoriser \(P\).
    \item Déterminer toutes ses racines dans \(\C\).
    \item Représenter ces racines dans le plan.
    \item Que peut-on dire du barycentre de ces points ?
\end{enumerate}

\begin{corrige}
\textbf{1) Factorisation}

On teste les racines évidentes parmi les entiers divisant \(2\) : \(\pm 1,\pm 2\).

On calcule
\[
P(1)=1-3+2=0.
\]
Donc \(1\) est racine. Ainsi
\[
P(z)=(z-1)Q(z).
\]

Effectuons la division :
\[
z^3-3z+2=(z-1)(z^2+z-2).
\]
Puis
\[
z^2+z-2=(z-1)(z+2).
\]
Donc
\[
\boxed{P(z)=(z-1)^2(z+2)}.
\]

\textbf{2) Racines}

Les racines sont donc
\[
z=1 \quad (\text{racine double})
\qquad \text{et} \qquad
z=-2.
\]

\textbf{3) Représentation dans le plan}

Ce sont les points de l'axe réel d'abscisses \(1\) et \(-2\). Le point d'affixe \(1\) compte deux fois si l'on tient compte de la multiplicité.

\textbf{4) Barycentre}

La somme des racines d'un polynôme unitaire de degré \(3\) vaut l'opposé du coefficient de \(z^2\), ici \(0\). Donc
\[
1+1+(-2)=0.
\]
Le barycentre des trois points pondérés par leur multiplicité a donc pour affixe
\[
\frac{1+1-2}{3}=0.
\]
Ainsi, leur barycentre est l'origine.

\[
\boxed{\text{Le barycentre des racines, comptées avec multiplicité, est l'origine.}}
\]
\end{corrige}

% ============================================================
\section{Exercice 20 \threeStars\ --- Minimum d'une somme de distances}
\difficulty{\threeStars}

Déterminer le minimum de
\[
|z-(1+\ii)|+|z-(1-\ii)|
\]
quand \(z\) parcourt \(\C\).

\begin{corrige}
Soit \(M\) le point d'affixe \(z\), \(A\) le point d'affixe \(1+\ii\), \(B\) le point d'affixe \(1-\ii\).

L'expression à minimiser est
\[
MA+MB.
\]

Par l'inégalité triangulaire, on sait que pour tout point \(M\),
\[
MA+MB \geq AB.
\]
Or
\[
AB=|(1+\ii)-(1-\ii)|=|2\ii|=2.
\]
Donc
\[
MA+MB\geq 2.
\]

Il reste à voir si cette borne est atteinte. L'égalité dans l'inégalité triangulaire est atteinte lorsque \(M\) est situé sur le segment \([AB]\).

Donc le minimum vaut \(2\), et il est atteint pour tout point \(M\) du segment reliant \(A\) à \(B\).

\textbf{Conclusion :}
\[
\boxed{\min_{z\in\C}\left(|z-(1+\ii)|+|z-(1-\ii)|\right)=2.}
\]

L'ensemble des points où ce minimum est atteint est le segment \([AB]\), c'est-à-dire le segment vertical d'équation \(x=1\) entre \(y=-1\) et \(y=1\).
\end{corrige}

% ============================================================
\section*{Conclusion générale}

Ces exercices couvrent plusieurs niveaux de maîtrise des nombres complexes :
\begin{itemize}[label=--]
    \item calcul algébrique ;
    \item résolution d'équations ;
    \item forme trigonométrique ;
    \item racines de l'unité ;
    \item géométrie complexe ;
    \item raisonnements de transition vers la prépa.
\end{itemize}

\textbf{Conseil de travail :}
\begin{enumerate}
    \item savoir refaire tous les calculs sans regarder le corrigé ;
    \item être capable d'expliquer chaque raisonnement à l'oral ;
    \item repérer les idées-type :
    \begin{itemize}
        \item poser \(z=x+\ii y\),
        \item utiliser le conjugué,
        \item traduire \(|z-a|\) comme une distance,
        \item traduire un quotient comme une rotation + homothétie,
        \item utiliser la forme trigonométrique pour les puissances.
    \end{itemize}
\end{enumerate}

\end{document}